% !TeX root = linear-codes.tex
% The binary matrix-vector product illustrated with marks and empty cells.
\begin{center}
\begin{tikzpicture}[x=2.3mm,y=2.3mm]
  \tikzset{binary grid line/.style={
    draw=black!55,
    line width=.22pt,
    line cap=butt
  }}
  \def\zeromarkhalfwidth{.36}
  \def\zeromarkcorner{.90pt}

  % Draw a binary grid from row-major data, with the top row first.
  % Arguments: lower-left x, lower-left y, rows, columns, data.
  % Ones are centered purple rounded squares that do not touch the cell walls.
  % The complete grid is drawn afterward as
  % rows+1 uninterrupted horizontal lines and columns+1 vertical lines.
  \newcommand{\drawbinarygrid}[5]{%
    \path[fill=white] (#1,#2) rectangle ++(#4,#3);
    \edef\binarydata{#5}
    \foreach \bit [count=\cell from 0] in \binarydata {%
      \pgfmathtruncatemacro{\gridrow}{int(\cell/#4)}
      \pgfmathtruncatemacro{\gridcol}{mod(\cell,#4)}
      \ifnum\bit=1\relax
        \path[fill=purple,rounded corners=\zeromarkcorner]
          ({#1+\gridcol+.5-\zeromarkhalfwidth},
           {#2+#3-\gridrow-.5-\zeromarkhalfwidth})
          rectangle
          ({#1+\gridcol+.5+\zeromarkhalfwidth},
           {#2+#3-\gridrow-.5+\zeromarkhalfwidth});
      \fi
    }
    \foreach \gridrow in {0,...,#3} {
      \draw[binary grid line]
        (#1,{#2+\gridrow}) -- ++(#4,0);
    }
    \foreach \gridcol in {0,...,#4} {
      \draw[binary grid line]
        ({#1+\gridcol},#2) -- ++(0,#3);
    }
  }

  \def\matrixrows{10}
  \def\matrixcols{18}
  \def\matrixleft{7}
  \def\matrixbottom{1}

  \def\xdata{0,1,0,0,1,0,0,0,0,0,0,0,0,1,0,1,0,0}
  \def\ydata{0,0,0,0,0,0,0,1,1,1}
  \def\matrixdata{
    1,1,0,1,0,0,1,1,0,0,1,1,0,1,1,0,0,1,
    0,0,1,0,1,1,1,0,1,1,1,0,1,0,0,1,1,0,
    0,0,1,1,0,1,0,1,1,0,0,1,1,0,1,0,1,1,
    1,1,0,0,1,0,1,1,0,1,1,0,0,1,0,1,0,0,
    1,1,1,0,0,1,0,0,1,1,0,1,1,0,0,1,1,0,
    0,0,1,1,1,0,1,1,1,0,1,0,0,1,1,0,0,1,
    1,1,1,0,0,0,0,1,0,1,1,0,1,0,0,1,1,0,
    0,0,0,1,1,1,1,0,1,0,1,1,0,0,1,0,0,1,
    0,1,0,1,1,0,1,0,1,1,0,1,0,1,1,0,1,1,
    1,1,1,0,0,0,1,1,0,1,0,0,1,1,0,1,0,0
  }

  % Derive the layout from the matrix dimensions.
  \pgfmathsetmacro{\matrixcenterx}{\matrixleft+\matrixcols/2}
  \pgfmathsetmacro{\matrixcentery}{\matrixbottom+\matrixrows/2}
  \pgfmathsetmacro{\matrixtop}{\matrixbottom+\matrixrows}
  \pgfmathsetmacro{\xgridbottom}{\matrixtop+1.5}
  \pgfmathsetmacro{\xgridtop}{\xgridbottom+1}
  \pgfmathsetmacro{\legendheight}{\matrixbottom-.45}

  % x is aligned with the matrix columns and separated by 1.5 cells.
  \drawbinarygrid{\matrixleft}{\xgridbottom}
    {1}{\matrixcols}{\xdata}
  \drawbinarygrid{0}{\matrixbottom}
    {\matrixrows}{1}{\ydata}
  \drawbinarygrid{\matrixleft}{\matrixbottom}
    {\matrixrows}{\matrixcols}{\matrixdata}

  \node[left=2.5pt] at (0,\matrixcentery) {$y$};
  \node at (3.8,\matrixcentery) {$=$};
  \node[above=1.5pt] at (\matrixcenterx,\xgridtop) {$x$};
\end{tikzpicture}

\vspace{2mm}
{\footnotesize The filled cells of $x$ select columns of $M$;
their sum is $y=Mx$.}
\end{center}
